\setchoices{label=(\Alph*)}

\begin{question}[index=1]
设函数 $f(x)=x-\ln x-ax^2$（$a\ne0$）有唯一零点，则 $a$ 的取值范围为
\begin{choices}[columns=4]
\item $(0,2)$
\item $(0,+\infty)$
\item $(2,+\infty)$
\item $(2,\eu)$
\end{choices}
\end{question}
\QuestionBreak

\begin{question}[index=2]
设 $z=z(x,y)$ 由方程 $\dfrac xz=f\!\left(\dfrac yz\right)$ 确定，其中 $f$ 为可微函数，则 $x\dfrac{\partial z}{\partial x}+y\dfrac{\partial z}{\partial y}=$
\begin{choices}[columns=4]
\item $0$
\item $x$
\item $y$
\item $z$
\end{choices}
\end{question}
\QuestionBreak

\begin{question}[index=3]
设 $L$ 为 $2x^2+3y^2=1$，取逆时针方向，则 $\displaystyle I=\oint_L\frac{x\dd y-y\dd x}{3x^2+4y^2}=$
\begin{choices}[columns=2]
\item $\dfrac\pi{\sqrt2}$
\item $\dfrac\pi2$
\item $1$
\item $\dfrac\pi{\sqrt3}$
\end{choices}
\end{question}
\QuestionBreak

\begin{question}[index=4]
设可导函数 $y=f(x)$ 在 $[a,b]$ 上严格单调递增，$g(y)$ 是 $f(x)$ 的反函数，则 $\displaystyle\int_a^bf(x)\dd x+\int_{f(a)}^{f(b)}g(y)\dd y=$
\begin{choices}[columns=4]
\item $f(b)-f(a)$
\item $bf(b)-af(a)$
\item $af(b)-bf(a)$
\item $af(a)-bf(b)$
\end{choices}
\end{question}
\QuestionBreak

\begin{question}[index=5]
设 $4$ 维实列向量 $\vect\beta_1,\vect\beta_2,\vect\beta_3$ 线性无关，非零实列向量 $\vect\alpha_1,\vect\alpha_2,\vect\alpha_3,\vect\alpha_4$ 与 $\vect\beta_1,\vect\beta_2,\vect\beta_3$ 均正交，记 $\vect A=(\vect\alpha_1,\vect\alpha_2,\vect\alpha_3,\vect\alpha_4)$，$\vect A^*$ 是 $\vect A$ 的伴随矩阵，则下列选项正确的是
\begin{choices}[columns=2]
\item 方程组 $\vect A^*\vect X=\vect0$ 的解均是方程组 $\vect A\vect X=\vect0$ 的解
\item 方程组 $\vect A^*\vect X=\vect0$ 与方程组 $\vect A\vect X=\vect0$ 同解
\item 方程组 $\vect A\vect X=\vect0$ 的解均是方程组 $\vect A^*\vect X=\vect0$ 的解
\item 方程组 $\vect A\vect X=\vect0$ 的解不是方程组 $\vect A^*\vect X=\vect0$ 的解
\end{choices}
\end{question}
\QuestionBreak

\begin{question}[index=6]
设 $\vect A$ 与 $\vect B$ 相似，$\lambda\in\mathbb R$，则
\begin{choices}[columns=2]
\item $\lambda\vect E-\vect A=\lambda\vect E-\vect B$
\item $\vect A$ 与 $\vect B$ 均与同一个对角矩阵相似
\item $\lambda\vect E-\vect A\sim\lambda\vect E-\vect B$
\item 对于 $\vect A$ 与 $\vect B$ 的同一个特征值，$\vect A$ 与 $\vect B$ 均有相同的特征向量
\end{choices}
\end{question}
\QuestionBreak

\begin{question}[index=7]
设 $\vect\alpha=(1,1)^{\mathrm T}$，$4$ 阶矩阵 $\vect A$ 相似于 $\vect B=\mat{\vect\alpha\vect\alpha^{\mathrm T}&\vect O\\\vect O&2\vect\alpha\vect\alpha^{\mathrm T}}$，$r_1=r(\vect A)$，$r_2=r(\vect A-\vect E)$，$r_3=r(\vect A-2\vect E)$，则
\begin{choices}[columns=4]
\item $r_1<r_2<r_3$
\item $r_2<r_3<r_1$
\item $r_3<r_1<r_2$
\item $r_1<r_3<r_2$
\end{choices}
\end{question}
\QuestionBreak

\begin{question}[index=8]
设随机变量 $X$ 与 $Y$ 相互独立，且 $X\dist P(\lambda_1)$，$Y\dist P(\lambda_2)$．
\par 在已知 $X+Y=n$（$n$ 为正整数）的条件下，$X$ 服从的分布为
\begin{choices}[columns=1]
\item $B\!\left(n,\dfrac{\lambda_2}{\lambda_1+\lambda_2}\right)$
\item $B\!\left(n,\dfrac{\lambda_1}{\lambda_1+\lambda_2}\right)$
\item $B\!\left(n,\dfrac{\lambda_1\lambda_2}{(\lambda_1+\lambda_2)^2}\right)$
\item $B\!\left(n,\dfrac{\lambda_1}{(\lambda_1+\lambda_2)^2}\right)$
\end{choices}
\end{question}
\QuestionBreak

\begin{question}[index=9]
设 $\widehat\theta$ 是参数 $\theta$ 的无偏估计量，则“$D(\widehat\theta)>0$”是“$\widehat\theta^2$ 不是 $\theta^2$ 的无偏估计量”的
\begin{choices}[columns=4]
\item 充分不必要条件
\item 必要不充分条件
\item 充分必要条件
\item 不充分也不必要条件
\end{choices}
\end{question}
\QuestionBreak

\begin{question}[index=10]
设 $X$ 与 $Y$ 分别表示 $100$ 次独立重复试验成功的次数与失败的次数．当 $D(X)$ 和 $D(Y)$ 取得最大值时，$\operatorname{Cov}(X,Y)=$
\begin{choices}[columns=4]
\item $5$
\item $25$
\item $-25$
\item $-5$
\end{choices}
\end{question}
